% fig-d1-series-map.tex --- D1, the Clebsch series map.
%
% Defines \ClebschSeriesMap{<node>}, where <node> is one of
% I, II, III, IV, V, or none.  The named node is drawn with a heavier rule
% and a light grey tint ("you are here"); every other node is plain.
%
% Current proved series geometry.  Paper IV remains a logically independent
% reconstruction branch; no arrow identifies its plane with the companion
% pencil.

\providecommand{\ClebschSeriesMap}[1]{%
\csfSetHighlight{#1}%
\begin{tikzpicture}[
  every node/.style={outer sep=0pt},
]

\node[\csfHL{I},text width=34mm] (pI) at (-5.1,2.45)
  {\textbf{Paper I}\\deep-hole syndrome\\conference companion};
\node[\csfHL{II},text width=34mm] (pII) at (-5.1,1.0)
  {\textbf{Paper II}\\quadratic trade\\chordal companion};
\node[\csfHL{III},text width=34mm] (pIII) at (-5.1,-0.45)
  {\textbf{Paper III}\\arithmetic shadows\\conference companion};

\node[\csfHL{V},text width=40mm] (pV) at (0,1.0)
  {\textbf{Paper V}\\marked companion correspondence\\selected \(L\) and \(q-1\)};
\node[csfcentre,text width=39mm] (carrier) at (5.1,1.0)
  {\textbf{ONE SIX-AXIS CARRIER}\\distinct conference and chordal
   singular shadows};

\node[\csfHL{IV},text width=39mm] (pIV) at (-2.65,-2.25)
  {\textbf{Paper IV}\\minimum-word pairs};
\node[csfcentre,text width=39mm] (plane) at (2.65,-2.25)
  {marked \(\PGtwo{13}\),\\conic, and polarity};

\draw[csfbiarr] (pI) -- ($(pV.west)+(0,5mm)$);
\draw[csfbiarr] (pII) -- (pV.west);
\draw[csfbiarr] (pIII) -- ($(pV.west)+(0,-5mm)$);
\draw[csfarr] (pV) -- (carrier);
\draw[csfarr] (pIV) -- (plane);

\end{tikzpicture}%
}
